\subsection{Airports}
\label{sec:3_3_ingestion_airports}

\par
I started by reading \texttt{airports.dat} from the data directory mounted
inside the Spark container at \texttt{/data/}. Since the file is in CSV format
and has no header row, I loaded it with the CSV connector, inferred the schema,
and assigned the expected column names (same as with Airlines) (\cref{fig:3-3-csv-read}):

\begin{figure}[H]
	\centering
	\includegraphics[width=0.9\textwidth]{3_3_csv_read.png}
	\caption{Loading \texttt{airports.dat} as a DataFrame and displaying the
		first rows}
	\label{fig:3-3-csv-read}
\end{figure}

\par
Then I defined and registered a UDF that determines the geographic quadrant of
each airport based on its latitude and longitude. Using this UDF, I added the
\texttt{hemisphere\_quadrant} column to the DataFrame (\cref{fig:3-3-udfs}):

\begin{figure}[H]
	\centering
	\includegraphics[width=0.9\textwidth]{3_3_udfs.png}
	\caption{Creating the \texttt{hemisphere\_quadrant} column with a UDF}
	\label{fig:3-3-udfs}
\end{figure}

\par
After that, I stored the resulting DataFrame in the \texttt{airports} table of
the Postgres database using Spark's JDBC connector. Once the write completed, I
read the table back from Postgres using the \texttt{partitionColumn},
\texttt{lowerBound}, \texttt{upperBound}, and \texttt{numPartitions} options,
and displayed the first five rows to verify that the ingestion was successful (\cref{fig:3-3-postgres-write}):

\begin{figure}[H]
	\centering
	\includegraphics[width=0.9\textwidth]{3_3_postgres_write.png}
	\caption{Writing the airports DataFrame to Postgres and verifying the stored
		data, using \texttt{partitionColumn}, \texttt{lowerBound},
		\texttt{upperBound} and \texttt{numPartitions} on the read}
	\label{fig:3-3-postgres-write}
\end{figure}

\par
By default, \texttt{spark.read.jdbc} opens one connection and uses a single
task to pull the entire table, even if I have more than one executor
available. Adding the following options, I can tell Spark to split the read
into several parallel range queries instead, each with its own connection:

\begin{itemize}
	\item \texttt{partitionColumn = "airport\_id"}: this is a numeric primary
	      key, which is good for Spark since it can split it into clean,
	      non-overlapping ranges without scanning the table first. These ids
	      are fairly spread out, so the partitions end up similar in size.
	\item \texttt{lowerBound = 1}, \texttt{upperBound = 14000}:
	      \texttt{airport\_id} actually runs from 1 to 14110, but these are
	      only used as an indicator for how wide the partitions should be for
	      Spark, they don't drop any rows outside that range, ids below 1 or
	      above 14000 still end up in the first or last partition.
	\item \texttt{numPartitions = 4}: mirrors the number of cores per worker
	      in the cluster (\cref{sec:2_5_docker_spark_cluster}) so the parallel reads don't compete for
	      the same resources.
\end{itemize}

\par
Instead of pulling all 7{,}700 rows through a single connection, this runs
four reads in parallel, each grabbing its own slice of \texttt{airport\_id}s.
That reduces the chance of the connection becoming a bottleneck, especially
useful since this table gets joined against \texttt{routes} later in
\cref{sec:3_4_ingestion_routes}.
